%%
%% Author: shoupu_wan
%% 7/14/18
%%

\chapter{Mindset of Dynamic programming} \label{ch:dp}

Beginners like me would be intimidated by dynamic programming.
It seems what DP requires is a power of magic that creates multiple things at once and I just don't
have enough hands to handle all that at once.
DP needs a split of me doing one part of the thing and a split of me doing tightly coupled another
while the two splits coordinate closely.

\section{Algorithm Refactor}\label{sec:algorithmRefactor}

\begin{defn}
    When a formula is too complicated, it may be refactored into several formalae, with auxiliary
    formulae which may make formulation of the problem simpler and the concept easier to understand.
    We call this strategy \emph{formula refactoring}.
\end{defn}

\begin{defn}
    When a dynamic programming formulation for a problem becomes too complicated, one may
    introduce a vector variable for each subproblem.
    Each component of the vector variable has a dedicated and meaningful purpose.
    The subproblem is to be solved to obtain the vector variable.
\end{defn}

\section{Two ways of dynamic programming}\label{sec:twoWaysOfDynamicProgramming}

The problem is asking for dynamic programming.
There are generally two ways of attacking dynamic programming.
Think of what's the optimal solution for subproblems and then grow the subproblem into the
full-sized problem;
think of what's the suboptimal solution for the same problem with certain
constraints and then seek to remove the constraints to obtain the optimal solution.
The first approach is 'problem-oriented' and the second is 'solution-oriented'.

There are a row of houses, each house can be painted with three colors red, blue and green
\label{Evernote:Paint House (I and II)}
\begin{itemize}
    \item Problem-oriented\label{Evernote:Split "bothearthandsaturnspin" into words word break wordbreak}
    \item \item Solution-oriented
\end{itemize}

\section{Case Study--Tiling a Rectangle} \label{sec:tiling-rectangle}
\label{Evernote:2x1 brick paving problem tiling}
2x1 Bricks covering a rectangular area.

\section{Two-sum Problem}

There is a greedy algorithm for solving this problems. But why does it work at all?

\section{Knapsack Problem}

The Knapsack problem is usually formulated as follows:
\begin{quote}
    \url{https://en.wikipedia.org/wiki/Knapsack_problem}
    Given a set of items, each with a weight and a value, determine the number of each item to
    include in a collection so that the total weight is less than or equal to a given limit and the
    total value is as large as possible.
\end{quote}

This problem is $NP$-hard.
But its limited form-integer weight and value-has a pseudo-polynomial solution.
A classical way to solve this problem is dynamic programming.
The question is ``Why the solution works at all?''
Why order of picking items affects neither the performance nor the result?
